\documentclass[aspectratio=169]{ctexbeamer}

\usepackage{amsmath}
\usepackage{booktabs}
\usepackage{graphicx}
\usepackage{listings}
\usepackage{tikz}
\usepackage{xcolor}
\usetikzlibrary{arrows.meta, positioning, shapes.geometric}

\usetheme{Berlin}
\usecolortheme{default}

\definecolor{mainblue}{RGB}{32, 83, 154}
\definecolor{softblue}{RGB}{230, 238, 250}
\definecolor{softgreen}{RGB}{232, 245, 233}
\definecolor{softorange}{RGB}{255, 243, 224}
\definecolor{codebg}{RGB}{248, 248, 246}
\definecolor{codekw}{RGB}{126, 63, 160}
\definecolor{codecomment}{RGB}{70, 130, 70}
\definecolor{codered}{RGB}{160, 60, 60}

\setbeamercolor{structure}{fg=mainblue}
\setbeamercolor{block title}{fg=white,bg=mainblue}
\setbeamercolor{block body}{bg=softblue}
\setbeamertemplate{navigation symbols}{}

\lstdefinelanguage{RustLike}{
    morekeywords={fn, let, mut, if, else, while, for, in, loop, break, continue, return, i32},
    sensitive=true,
    morecomment=[l]{//},
    morecomment=[s]{/*}{*/},
    morestring=[b]"
}

\lstdefinestyle{code}{
    language=RustLike,
    backgroundcolor=\color{codebg},
    basicstyle=\ttfamily\scriptsize,
    keywordstyle=\color{codekw}\bfseries,
    commentstyle=\color{codecomment},
    stringstyle=\color{codered},
    showstringspaces=false,
    breaklines=true,
    columns=fullflexible,
    keepspaces=true,
    frame=single,
    rulecolor=\color{black!18},
    xleftmargin=0.5em,
    xrightmargin=0.5em
}

\lstdefinestyle{tinycode}{
    style=code,
    basicstyle=\ttfamily\tiny
}

\newcommand{\speaker}[2]{%
    \vspace{-0.4em}
    \begin{flushright}
        \scriptsize\color{mainblue}\textbf{讲者：#1}\quad 建议用时：#2
    \end{flushright}
    \vspace{-0.7em}
}

\title[编译原理大作业 1 \& 2]{编译原理大作业 1 \& 2}
\subtitle{类 Rust 语言的词法/语法分析、语义分析与四元式中间代码生成}
\author[李堂辉, 高胜寒, 冉子易]{2351078 李堂辉 \\ 2351837 高胜寒 \\ 2352197 冉子易\\
\small （按学号为序）}
\institute{同济大学}
\date{2026年6月16日}

\begin{document}

\begin{frame}
    \titlepage
\end{frame}

\begin{frame}{汇报主线：两个作业是一条编译流程}
    \speaker{冉子易}{25 秒}
    \begin{center}
    \resizebox{0.95\linewidth}{!}{
    \begin{tikzpicture}[
        node distance=1.5cm,
        box/.style={rectangle, rounded corners, draw=mainblue, thick, minimum width=2.45cm, minimum height=0.85cm, align=center, fill=softblue},
        box2/.style={rectangle, rounded corners, draw=mainblue, thick, minimum width=2.45cm, minimum height=0.85cm, align=center, fill=softgreen},
        arrow/.style={-{Latex[length=3mm]}, thick, draw=mainblue}
    ]
        \node[box] (src) {类 Rust\\源程序};
        \node[box, right=of src] (tok) {词法分析\\Token 序列};
        \node[box, right=of tok] (ast) {语法分析\\AST};
        \node[box2, right=of ast] (sem) {语义分析\\错误检查};
        \node[box2, right=of sem] (ir) {中间代码\\四元式 IR};
        \draw[arrow] (src) -- (tok);
        \draw[arrow] (tok) -- (ast);
        \draw[arrow] (ast) -- (sem);
        \draw[arrow] (sem) -- (ir);
        \node[below=0.35cm of tok, align=center, text=mainblue] {\small 大作业 1};
        \node[below=0.35cm of sem, align=center, text=mainblue] {\small 大作业 2};
    \end{tikzpicture}}
    \end{center}
    \vspace{-0.2em}
    \begin{table}
    \centering
    \small
    \begin{tabular}{lll}
        \toprule
        \textbf{阶段} & \textbf{核心问题} & \textbf{产物} \\
        \midrule
        大作业 1 & 源码是否符合词法/语法规则 & Token、AST \\
        大作业 2 & AST 是否语义正确，如何线性化 & 语义诊断、四元式 IR \\
        \bottomrule
    \end{tabular}
    \end{table}
\end{frame}

\begin{frame}{8 分钟分工}
    \speaker{冉子易}{20 秒}
    \begin{table}
    \centering
    \small
    \begin{tabular}{p{2.2cm}p{2.0cm}p{7.6cm}}
        \toprule
        \textbf{讲者} & \textbf{时间} & \textbf{技术内容} \\
        \midrule
        冉子易 & 0:00--2:00 & 两个作业关系；大作业 1 的 Token、Lexer、递归下降 Parser、AST。\\
        高胜寒 & 2:00--4:00 & 大作业 2 的语义分析：符号表栈、作用域、类型检查、可变性检查、错误收集。\\
        李堂辉 & 4:00--6:00 & 四元式 IR 设计、表达式翻译、控制流翻译、测试结果和总结。\\
        全员 & 6:00--8:00 & 按模块问答：前端 / 语义 / IR 分别回答。\\
        \bottomrule
    \end{tabular}
    \end{table}
\end{frame}

\begin{frame}{核心文件地图}
    \speaker{冉子易}{20 秒}
    \begin{table}
    \centering
    \small
    \begin{tabular}{p{2.5cm}p{3.0cm}p{5.7cm}}
        \toprule
        \textbf{位置} & \textbf{文件} & \textbf{职责} \\
        \midrule
        大作业 1 & \texttt{token.rs} & Token 类型、显示格式、类别划分 \\
        大作业 1 & \texttt{lexer.rs} & 字符扫描、关键字识别、注释跳过、错误定位 \\
        大作业 1/2 & \texttt{parser.rs, ast.rs} & 递归下降解析、AST 节点定义 \\
        大作业 2 & \texttt{compiler.rs} & 符号表、语义检查、AST 到 IR 的翻译 \\
        大作业 2 & \texttt{ir.rs} & \texttt{Op / Arg / Quadruple / IRProgram} \\
        \bottomrule
    \end{tabular}
    \end{table}
\end{frame}

\section{大作业 1：词法与语法分析}

\begin{frame}{大作业 1 的目标与范围}
    \speaker{冉子易}{20 秒}
    \begin{columns}[T]
    \begin{column}{0.48\textwidth}
        \begin{block}{输入}
            类 Rust 语言子集，包括：
            \begin{itemize}
                \item 函数定义、参数、返回类型
                \item 变量声明、赋值、表达式
                \item \texttt{if / while / for / loop}
                \item 数组、元组、引用与解引用
            \end{itemize}
        \end{block}
    \end{column}
    \begin{column}{0.48\textwidth}
        \begin{block}{输出}
            \begin{itemize}
                \item Token 序列：带类别、文本、行列号
                \item AST：把程序组织成树形结构
                \item 词法/语法错误定位
            \end{itemize}
        \end{block}
    \end{column}
    \end{columns}
    \vspace{0.5em}
    \textbf{对应文件：}\texttt{token.rs, lexer.rs, parser.rs, ast.rs}
\end{frame}

\begin{frame}{词法分析：从字符流到 Token}
    \speaker{冉子易}{25 秒}
    \begin{columns}[T]
    \begin{column}{0.56\textwidth}
        \begin{itemize}
            \item 手写扫描器，逐字符推进，并维护当前位置。
            \item 先按首字符分类，再对二字符符号做前瞻。
            \item 标识符先整体读入，再查关键字表。
            \item 注释、空白符不进入最终 Token 序列。
        \end{itemize}
    \end{column}
    \begin{column}{0.4\textwidth}
        \begin{block}{典型区分}
            \small
            \begin{tabular}{ll}
                \texttt{=} & 赋值 \\
                \texttt{==} & 相等比较 \\
                \texttt{.} & 元组字段访问 \\
                \texttt{..} & 区间表达式 \\
                \texttt{-} & 减号 \\
                \texttt{->} & 函数返回类型 \\
            \end{tabular}
        \end{block}
    \end{column}
    \end{columns}
\end{frame}

\begin{frame}[fragile]{Token 设计：枚举直接表达语法类别}
    \speaker{冉子易}{25 秒}
    \begin{lstlisting}[style=code]
pub enum TokenKind {
    Fn, Let, Mut, If, Else, While, For, In, Loop,
    Break, Continue, Return, I32,
    IntLiteral(i64),
    Ident(String),
    Assign, Plus, Minus, Star, Slash,
    Eq, Ne, Lt, Le, Gt, Ge,
    Ampersand,
    LParen, RParen, LBrace, RBrace, LBracket, RBracket,
    Semicolon, Colon, Comma,
    Arrow, Dot, DotDot,
    Eof,
}
    \end{lstlisting}
    \vspace{-0.4em}
    \textbf{设计点：}整数和标识符直接携带值，后续 Parser 不需要再查原始字符串。
\end{frame}

\begin{frame}{语法分析：递归下降 + 优先级分层}
    \speaker{冉子易}{30 秒}
    \begin{columns}[T]
    \begin{column}{0.52\textwidth}
        \begin{block}{递归下降}
            每个非终结符对应一个解析函数，按当前 Token 选择产生式，成功后推进 Token 指针。
        \end{block}
        \begin{itemize}
            \item \texttt{parse\_program}
            \item \texttt{parse\_function\_decl}
            \item \texttt{parse\_statement}
            \item \texttt{parse\_expression}
        \end{itemize}
    \end{column}
    \begin{column}{0.44\textwidth}
        \begin{block}{表达式优先级}
            \small
            \begin{tabular}{ll}
                低 & 区间 \texttt{..} \\
                   & 比较 \texttt{< <= == !=} \\
                   & 加减 \texttt{+ -} \\
                   & 乘除 \texttt{* /} \\
                   & 一元 \texttt{*x, \&x} \\
                高 & 调用 / 索引 / 元组访问 \\
            \end{tabular}
        \end{block}
    \end{column}
    \end{columns}
\end{frame}

\begin{frame}[fragile]{AST：把语法结构保留下来}
    \speaker{冉子易}{20 秒}
    \begin{columns}[T]
    \begin{column}{0.49\textwidth}
        \begin{lstlisting}[style=tinycode]
fn add(mut a:i32, b:i32) -> i32 {
    let mut c = a + b;
    return c;
}
        \end{lstlisting}
    \end{column}
    \begin{column}{0.47\textwidth}
        \begin{itemize}
            \item 顶层是 \texttt{Program}
            \item 函数是 \texttt{FunctionDecl}
            \item 语句是 \texttt{Statement}
            \item 表达式是 \texttt{Expr}
            \item 二元运算保留左右子树和操作符
        \end{itemize}
    \end{column}
    \end{columns}
    \begin{block}{为什么需要 AST}
        后续语义分析和 IR 生成不再处理字符或 Token，而是遍历结构化的 AST。
    \end{block}
\end{frame}

\begin{frame}{大作业 1 数据流}
    \speaker{冉子易}{20 秒}
    \begin{center}
    \resizebox{0.92\linewidth}{!}{
    \begin{tikzpicture}[
        node distance=1.0cm,
        box/.style={rectangle, rounded corners, draw=mainblue, thick, minimum width=2.9cm, minimum height=0.85cm, align=center, fill=softblue},
        arr/.style={-{Latex[length=2.8mm]}, thick, draw=mainblue}
    ]
        \node[box] (src) {\texttt{fn f(a:i32)\{...\}}\\源代码};
        \node[box, right=of src] (chars) {字符流\\行列号};
        \node[box, right=of chars] (tokens) {Token 序列\\类别 + 字面值};
        \node[box, right=of tokens] (ast) {AST\\结构化语法树};
        \draw[arr] (src) -- (chars);
        \draw[arr] (chars) -- node[above]{Lexer} (tokens);
        \draw[arr] (tokens) -- node[above]{Parser} (ast);
    \end{tikzpicture}}
    \end{center}
    \begin{table}
    \centering
    \small
    \begin{tabular}{lll}
        \toprule
        \textbf{错误类型} & \textbf{发现位置} & \textbf{例子} \\
        \midrule
        非法字符 & Lexer & \texttt{\$} \\
        Token 不匹配 & Parser & 缺少 \texttt{;} 或 \texttt{\}} \\
        表达式结构错误 & Parser & 括号不匹配、操作数缺失 \\
        \bottomrule
    \end{tabular}
    \end{table}
\end{frame}

\section{大作业 2：语义分析}

\begin{frame}{大作业 2 在大作业 1 基础上的扩展}
    \speaker{高胜寒}{20 秒}
    \begin{center}
    \resizebox{0.9\linewidth}{!}{
    \begin{tikzpicture}[
        node distance=1.25cm,
        box/.style={rectangle, rounded corners, draw=mainblue, thick, minimum width=2.7cm, minimum height=0.75cm, align=center, fill=softblue},
        newbox/.style={rectangle, rounded corners, draw=mainblue, thick, minimum width=2.7cm, minimum height=0.75cm, align=center, fill=softorange},
        arrow/.style={-{Latex[length=2.8mm]}, thick, draw=mainblue}
    ]
        \node[box] (lexer) {\texttt{lexer.rs}\\Tokenize};
        \node[box, right=of lexer] (parser) {\texttt{parser.rs}\\Parse AST};
        \node[newbox, right=of parser] (compiler) {\texttt{compiler.rs}\\语义 + 生成};
        \node[newbox, right=of compiler] (ir) {\texttt{ir.rs}\\四元式};
        \draw[arrow] (lexer) -- (parser);
        \draw[arrow] (parser) -- (compiler);
        \draw[arrow] (compiler) -- (ir);
    \end{tikzpicture}}
    \end{center}
    \vspace{-0.2em}
    \begin{table}
    \centering
    \small
    \begin{tabular}{lll}
        \toprule
        \textbf{输入} & \textbf{新增状态} & \textbf{输出} \\
        \midrule
        AST & 符号表栈、循环栈、返回类型 & 语义错误列表或 IR \\
        \bottomrule
    \end{tabular}
    \end{table}
\end{frame}

\begin{frame}[fragile]{符号表：用作用域栈模拟变量生命周期}
    \speaker{高胜寒}{30 秒}
    \begin{columns}[T]
    \begin{column}{0.48\textwidth}
        \begin{lstlisting}[style=tinycode]
struct SymbolInfo {
    mutable: bool,
    ty: Type,
    arg: Arg,
}

struct Compiler {
    scopes: Vec<HashMap<String, SymbolInfo>>,
    errors: Vec<String>,
}
        \end{lstlisting}
    \end{column}
    \begin{column}{0.48\textwidth}
        \begin{itemize}
            \item 进入函数、代码块、表达式块时 \texttt{push\_scope}
            \item 退出时 \texttt{pop\_scope}
            \item 查询变量时从内层向外层倒序查找
            \item 支持变量重影，符合类 Rust 的直觉
        \end{itemize}
    \end{column}
    \end{columns}
\end{frame}

\begin{frame}{作用域查询：从内到外}
    \speaker{高胜寒}{25 秒}
    \begin{columns}[T]
    \begin{column}{0.43\textwidth}
        \begin{center}
        \resizebox{0.95\linewidth}{!}{
        \begin{tikzpicture}[
            scope/.style={rectangle, rounded corners, draw=mainblue, thick, minimum width=3.4cm, minimum height=0.85cm, align=center},
            arr/.style={-{Latex[length=2.8mm]}, thick, draw=mainblue}
        ]
            \node[scope, fill=softorange] (inner) {当前块\\\texttt{x: i32}};
            \node[scope, fill=softgreen, below=0.3cm of inner] (func) {函数作用域\\\texttt{a, b}};
            \node[scope, fill=softblue, below=0.3cm of func] (global) {全局作用域\\函数声明};
            \draw[arr] (inner.east) to[bend left=13] node[right]{查找} (func.east);
            \draw[arr] (func.east) to[bend left=13] (global.east);
        \end{tikzpicture}}
        \end{center}
    \end{column}
    \begin{column}{0.53\textwidth}
        \scriptsize
        \setlength{\tabcolsep}{3pt}
        \begin{tabular}{p{1.45cm}p{2.35cm}p{2.35cm}}
            \toprule
            \textbf{操作} & \textbf{动作} & \textbf{目的} \\
            \midrule
            进入块 & \texttt{push\_scope()} & 隔离局部变量 \\
            声明变量 & 插入当前表 & 记录类型/可变性 \\
            使用变量 & 反向查找 & 支持遮蔽和外层引用 \\
            退出块 & \texttt{pop\_scope()} & 结束生命周期 \\
            \bottomrule
        \end{tabular}
    \end{column}
    \end{columns}
\end{frame}

\begin{frame}{语义检查覆盖点}
    \speaker{高胜寒}{30 秒}
    \begin{table}
    \centering
    \small
    \begin{tabular}{p{3.1cm}p{4.3cm}p{4.2cm}}
        \toprule
        \textbf{错误类型} & \textbf{触发场景} & \textbf{检查位置} \\
        \midrule
        未声明变量 & \texttt{a = 1;} 但没有 \texttt{let a} & \texttt{lookup\_variable} \\
        不可变变量赋值 & \texttt{let a = 1; a = 2;} & \texttt{compile\_lvalue} \\
        赋值类型不匹配 & 左值类型与右值类型不同 & \texttt{compile\_statement} \\
        运算类型不匹配 & \texttt{i32 + tuple} & \texttt{compile\_expression} \\
        返回类型不一致 & \texttt{-> i32} 但 \texttt{return;} & \texttt{Statement::Return} \\
        非法 break/continue & 不在循环体内使用 & \texttt{loop\_stack} \\
        \bottomrule
    \end{tabular}
    \end{table}
\end{frame}

\begin{frame}[fragile]{语义错误示例}
    \speaker{高胜寒}{25 秒}
    \begin{columns}[T]
    \begin{column}{0.48\textwidth}
        \begin{lstlisting}[style=tinycode]
fn test_immutable_assign() {
    let a: i32 = 1;
    a = 2;
}

fn test_break_outside_loop() {
    break;
    continue;
}
        \end{lstlisting}
    \end{column}
    \begin{column}{0.48\textwidth}
        \begin{block}{检查逻辑}
            \begin{itemize}
                \item 赋值左侧必须是合法左值。
                \item 左值如果来自不可变变量，则报错。
                \item \texttt{break/continue} 必须能在 \texttt{loop\_stack} 中找到循环上下文。
            \end{itemize}
        \end{block}
    \end{column}
    \end{columns}
\end{frame}

\section{大作业 2：四元式中间代码}

\begin{frame}{第三部分：四元式 IR}
    \speaker{李堂辉}{15 秒}
    \begin{center}
    \resizebox{0.88\linewidth}{!}{
    \begin{tikzpicture}[
        box/.style={rectangle, rounded corners, draw=mainblue, thick, minimum width=3.0cm, minimum height=0.8cm, align=center, fill=softblue},
        irbox/.style={rectangle, rounded corners, draw=mainblue, thick, minimum width=3.0cm, minimum height=0.8cm, align=center, fill=softgreen},
        arr/.style={-{Latex[length=2.8mm]}, thick, draw=mainblue}
    ]
        \node[box] (ast) {AST\\结构化语法};
        \node[box, right=1.2cm of ast] (sem) {语义通过\\类型/作用域有效};
        \node[irbox, right=1.2cm of sem] (ir) {四元式 IR\\线性指令序列};
        \draw[arr] (ast) -- (sem);
        \draw[arr] (sem) -- (ir);
        \node[below=0.45cm of ir, align=center] {\texttt{(op, arg1, arg2, result)}};
    \end{tikzpicture}}
    \end{center}
    \begin{table}
    \centering
    \small
    \begin{tabular}{lll}
        \toprule
        \textbf{语法结构} & \textbf{IR 处理} & \textbf{关键状态} \\
        \midrule
        表达式 & 临时变量承接中间结果 & \texttt{temp\_count} \\
        分支/循环 & 标签 + 条件/无条件跳转 & \texttt{label\_count} \\
        break/continue & 跳转到循环边界标签 & \texttt{loop\_stack} \\
        \bottomrule
    \end{tabular}
    \end{table}
\end{frame}

\begin{frame}[fragile]{四元式数据结构}
    \speaker{李堂辉}{30 秒}
    \begin{columns}[T]
    \begin{column}{0.48\textwidth}
        \begin{lstlisting}[style=tinycode]
pub enum Op {
    Add, Sub, Mul, Div,
    Eq, Ne, Lt, Le, Gt, Ge,
    Assign,
    Jump, JumpIfFalse, Label,
    Param, Arg, Call, Return,
    Ref, Deref, DerefAssign,
    LoadArray, StoreArray,
}
        \end{lstlisting}
    \end{column}
    \begin{column}{0.48\textwidth}
        \begin{lstlisting}[style=tinycode]
pub enum Arg {
    Empty,
    Int(i64),
    Var(String),
    Temp(usize),
    Label(usize),
}

pub struct Quadruple {
    op: Op,
    arg1: Arg,
    arg2: Arg,
    result: Arg,
}
        \end{lstlisting}
    \end{column}
    \end{columns}
    \textbf{设计点：}表达式结果统一放进临时变量 \texttt{t0, t1, ...}，控制流目标统一用标签 \texttt{L0, L1, ...}。
\end{frame}

\begin{frame}{IR 指令族}
    \speaker{李堂辉}{20 秒}
    \begin{table}
    \centering
    \scriptsize
    \begin{tabular}{p{2.4cm}p{4.4cm}p{4.7cm}}
        \toprule
        \textbf{类别} & \textbf{指令} & \textbf{用途} \\
        \midrule
        算术/比较 & \texttt{ADD SUB MUL DIV}\\ \texttt{EQ NE LT LE GT GE} & 表达式计算，结果写入临时变量 \\
        数据移动 & \texttt{ASSIGN REF DEREF}\\ \texttt{DEREF\_ASSIGN LOAD\_ARRAY} & 变量赋值、引用、数组读取 \\
        控制流 & \texttt{LABEL JUMP}\\ \texttt{JUMP\_IF\_FALSE} & 表示分支、循环、\texttt{break/continue} \\
        函数 & \texttt{PARAM ARG CALL RETURN} & 参数、调用和返回 \\
        \bottomrule
    \end{tabular}
    \end{table}
\end{frame}

\begin{frame}[fragile]{表达式翻译：递归生成临时变量}
    \speaker{李堂辉}{30 秒}
    \begin{columns}[T]
    \begin{column}{0.4\textwidth}
        \begin{lstlisting}[style=code]
let mut c = a + b * 2;
        \end{lstlisting}
    \end{column}
    \begin{column}{0.56\textwidth}
        \begin{lstlisting}[style=code]
(MUL,    b,   2,   t0)
(ADD,    a,  t0,   t1)
(ASSIGN, t1,  -,    c)
        \end{lstlisting}
    \end{column}
    \end{columns}
    \vspace{0.2em}
    \begin{center}
    \small
    \begin{tabular}{ccl}
        \texttt{b * 2} & $\Rightarrow$ & \texttt{t0} \\
        \texttt{a + t0} & $\Rightarrow$ & \texttt{t1} \\
        \texttt{t1} & $\Rightarrow$ & \texttt{c} \\
    \end{tabular}
    \end{center}
\end{frame}

\begin{frame}[fragile]{if 翻译：条件跳转 + 两个标签}
    \speaker{李堂辉}{25 秒}
    \begin{columns}[T]
    \begin{column}{0.4\textwidth}
        \begin{lstlisting}[style=tinycode]
if a > 0 {
    return 1;
} else {
    return 0;
}
        \end{lstlisting}
    \end{column}
    \begin{column}{0.56\textwidth}
        \begin{lstlisting}[style=tinycode]
(GT,             a, 0, t0)
(JUMP_IF_FALSE, t0, -, L0)
(RETURN,         1, -, -)
(JUMP,           -, -, L1)
L0:
(RETURN,         0, -, -)
L1:
        \end{lstlisting}
    \end{column}
    \end{columns}
    \textbf{转换形态：}\texttt{条件失败跳转} $\rightarrow$ \texttt{then} $\rightarrow$ \texttt{跳过 else} $\rightarrow$ \texttt{else 标签}。
\end{frame}

\begin{frame}[fragile]{while/loop 翻译：循环栈管理 break 和 continue}
    \speaker{李堂辉}{35 秒}
    \begin{lstlisting}[style=tinycode]
let start_label = self.new_label();
let end_label = self.new_label();

emit(Label,       -,        -, Label(start_label));
let cond_arg = compile_expression(cond);
emit(JumpIfFalse, cond_arg, -, Label(end_label));

loop_stack.push((start_label, end_label, None));
compile_block(body);
loop_stack.pop();

emit(Jump,        -,        -, Label(start_label));
emit(Label,       -,        -, Label(end_label));
    \end{lstlisting}
    \vspace{-0.4em}
    \begin{itemize}
        \item \texttt{continue} 跳到 \texttt{start\_label}，重新判断循环条件。
        \item \texttt{break} 跳到 \texttt{end\_label}，离开循环。
    \end{itemize}
\end{frame}

\begin{frame}[fragile]{函数调用翻译：参数先入 IR}
    \speaker{李堂辉}{25 秒}
    \begin{columns}[T]
    \begin{column}{0.38\textwidth}
        \begin{lstlisting}[style=code]
foo(a, 1 + b)
        \end{lstlisting}
    \end{column}
    \begin{column}{0.58\textwidth}
        \begin{lstlisting}[style=code]
(ARG,   a,   -, -)
(ADD,   1,   b, t0)
(ARG,  t0,   -, -)
(CALL, foo,  2, t1)
        \end{lstlisting}
    \end{column}
    \end{columns}
    \vspace{0.2em}
    \begin{table}
    \centering
    \small
    \begin{tabular}{lll}
        \toprule
        \textbf{字段} & \textbf{含义} & \textbf{示例} \\
        \midrule
        \texttt{arg1} & 函数名 & \texttt{foo} \\
        \texttt{arg2} & 参数个数 & \texttt{2} \\
        \texttt{result} & 返回临时变量 & \texttt{t1} \\
        \bottomrule
    \end{tabular}
    \end{table}
\end{frame}

\begin{frame}{测试与结果}
    \speaker{李堂辉}{25 秒}
    \begin{columns}[T]
    \begin{column}{0.48\textwidth}
        \begin{block}{大作业 1}
            \begin{itemize}
                \item 内置测试覆盖函数定义、声明、表达式、控制流等规则。
                \item 输出 Token 序列与 AST。
                \item 错误测试覆盖非法字符、语法错误等。
            \end{itemize}
        \end{block}
    \end{column}
    \begin{column}{0.48\textwidth}
        \begin{block}{大作业 2}
            \begin{itemize}
                \item 功能测试生成四元式。
                \item 语义错误测试覆盖未声明变量、类型不匹配、不可变赋值等。
                \item 两个 Rust 工程均可通过 \texttt{cargo check}。
            \end{itemize}
        \end{block}
    \end{column}
    \end{columns}
\end{frame}

\begin{frame}{问答分派}
    \speaker{李堂辉}{15 秒}
    \begin{table}
    \centering
    \small
    \begin{tabular}{p{3.0cm}p{4.2cm}p{4.5cm}}
        \toprule
        \textbf{问题方向} & \textbf{负责回答} & \textbf{关键词} \\
        \midrule
        词法/语法/AST & 冉子易 & Token、递归下降、优先级、AST 节点 \\
        语义分析 & 高胜寒 & 符号表、作用域、类型、可变性 \\
        IR/控制流/测试 & 李堂辉 & 四元式、临时变量、标签、循环栈 \\
        \bottomrule
    \end{tabular}
    \end{table}
\end{frame}

\begin{frame}{总结}
    \speaker{李堂辉}{20 秒}
    \begin{block}{最终完成的编译流程}
        \centering
        源码 $\rightarrow$ Token $\rightarrow$ AST $\rightarrow$ 语义检查 $\rightarrow$ 四元式 IR
    \end{block}
    \begin{itemize}
        \item 大作业 1 完成前端基础：词法分析、递归下降语法分析和 AST 构造。
        \item 大作业 2 在 AST 上继续做静态语义检查，并生成四元式中间代码。
        \item 技术重点在三个衔接：Token 到 AST、AST 到语义信息、结构化控制流到线性 IR。
    \end{itemize}
    \vfill
    \centering
    \Large 谢谢老师，欢迎提问
\end{frame}

\appendix
\section{问答备用}

\begin{frame}{备用：如果问“大作业 1 和 2 有什么区别？”}
    \begin{table}
    \centering
    \small
    \begin{tabular}{llll}
        \toprule
        \textbf{项目} & \textbf{输入} & \textbf{核心状态} & \textbf{输出} \\
        \midrule
        大作业 1 & 源代码字符串 & Token 指针、AST 构造器 & Token + AST \\
        大作业 2 & AST & 符号表栈、循环栈、临时变量计数 & 语义错误或四元式 IR \\
        \bottomrule
    \end{tabular}
    \end{table}
\end{frame}

\begin{frame}{备用：为什么选择递归下降？}
    \begin{itemize}
        \item 课程给定的类 Rust 子集规模适中，手写递归下降足够清晰。
        \item 每类语法结构对应一个函数，便于定位错误和展示实现。
        \item 表达式优先级可以通过多层函数自然表达。
        \item 不依赖 parser generator，项目文件更少，汇报时更容易解释。
    \end{itemize}
\end{frame}

\begin{frame}{备用：为什么需要语义分析？}
    \begin{itemize}
        \item 语法只能判断结构是否匹配，不能判断名字是否声明、类型是否一致。
        \item 例如 \texttt{a = 1;} 在语法上像赋值语句，但如果 \texttt{a} 没有声明，语义上是错的。
        \item 例如 \texttt{let a = 1; a = 2;} 语法没问题，但违反不可变变量不能二次赋值。
        \item 所以语义分析要维护符号表、类型信息和控制流上下文。
    \end{itemize}
\end{frame}

\begin{frame}{备用：四元式有什么好处？}
    \begin{itemize}
        \item 统一表达赋值、算术、比较、跳转、函数调用等操作。
        \item 复杂表达式会被拆成若干条三地址形式，便于后续优化。
        \item 控制流被压平成标签和跳转，后续可以进一步映射到汇编或虚拟机指令。
        \item 对课程项目来说，四元式比直接生成机器码更适合展示编译中间阶段。
    \end{itemize}
\end{frame}

\begin{frame}{备用：现有限制}
    \begin{itemize}
        \item 项目重点是课程要求的前端、语义检查和 IR 展示，没有继续生成真实机器码。
        \item 部分复杂类型和函数签名检查采用简化实现，例如函数调用默认返回 \texttt{i32}。
        \item 数组写入、元组访问等 IR 仍可进一步扩展为更完整的地址计算。
        \item 这些限制不影响主线：词法/语法分析、语义检查、基础四元式生成已经贯通。
    \end{itemize}
\end{frame}

\end{document}
